% META: {"id": "1NSI-RES-EVAL-B-corrige", "chapitre": "1NSI-RESEAUX", "type_objet": "corrige_evaluation", "evaluation_ref": "1NSI-RES-EVAL-B", "capacites": ["P-ARCH-02A", "P-ARCH-02B", "P-ARCH-02C", "P-ARCH-04A", "P-ARCH-04B"], "mode_creation": "ex_nihilo", "sources_inspiration": [], "status": "verified"}
% BEGIN-VERIFY
% def decouper_en_paquets(message, taille_max):
%     nb_paquets = (len(message) + taille_max - 1) // taille_max
%     paquets = []
%     for i in range(nb_paquets):
%         donnees = message[i * taille_max:(i + 1) * taille_max]
%         paquets.append({"numero": i, "total": nb_paquets, "donnees": donnees})
%     return paquets
% paquets = decouper_en_paquets("PAQUETSRESEAU", 5)
% assert len(paquets) == 3
% reseau = {"X": ["Y"], "Y": ["X", "Z"], "Z": ["Y"]}
% def peuvent_communiquer(reseau, depart, arrivee, visites=None):
%     if visites is None:
%         visites = set()
%     if depart == arrivee:
%         return True
%     visites.add(depart)
%     for v in reseau.get(depart, []):
%         if v not in visites and peuvent_communiquer(reseau, v, arrivee, visites):
%             return True
%     return False
% assert peuvent_communiquer(reseau, "X", "Z") is True
% END-VERIFY
\begin{corrige}{1NSI-RES-EVAL-B-EX1}
\textbf{1.} $3$ paquets sont créés : \lstinline{"PAQUE"} (numéro 0), \lstinline{"TSRES"}
(numéro 1), \lstinline{"EAU"} (numéro 2).

\textbf{2.} Tentative $1$ : envoi de \lstinline{"M"} avec bit $0$ — réussit. Tentative
$2$ : envoi de \lstinline{"N"} avec bit $1$ — échoue. Tentative $3$ : retransmission de
\lstinline{"N"} avec bit $1$ — réussit. Tentative $4$ : envoi de \lstinline{"O"} avec bit
$0$ — réussit. Messages reçus : \lstinline{[(0, "M"), (1, "N"), (0, "O")]}.

\textbf{3.} Les paquets sont envoyés \textbf{un par un} : l'émetteur attend toujours l'ACK
avant d'envoyer le suivant. Un seul bit suffit donc à distinguer \og un nouveau paquet \fg{}
(bit différent du précédent) d'\og une retransmission \fg{} (même bit que le paquet
précédent, non encore confirmé).
\end{corrige}

\begin{corrige}{1NSI-RES-EVAL-B-EX2}
\textbf{1.} \lstinline{peuvent_communiquer(reseau, "X", "Z")} renvoie \lstinline{True} :
il existe un chemin \lstinline{X -> Y -> Z}.

\textbf{2.}
\begin{python}
def decider_ventilation(co2, seuil=800):
    return co2 > seuil
\end{python}

\textbf{3.} \lstinline{decider_ventilation(1000)} vaut \lstinline{True} (air vicié, il
faut ventiler) ; \lstinline{decider_ventilation(500)} vaut \lstinline{False} (air
correct). Le \textbf{capteur} est le capteur de CO2 ; l'\textbf{actionneur} est le système
de ventilation.
\end{corrige}

\begin{corrige}{1NSI-RES-EVAL-B-EX3}
\textbf{1.}
\begin{python}
etat = {"chauffage": False}

def decider_chauffage(temperature, seuil_bas=24, seuil_haut=26):
    if etat["chauffage"]:
        if temperature > seuil_haut:
            etat["chauffage"] = False
    else:
        if temperature < seuil_bas:
            etat["chauffage"] = True
    return etat["chauffage"]
\end{python}

\textbf{2.} \lstinline{[False, True, True, False, False]}. À la troisième mesure
($25\,^\circ\text{C}$) le chauffage est déjà allumé (mesure précédente : $23 < 24$) et
$25$ ne dépasse pas $26$ : il reste allumé. La réponse dépend de l'état mémorisé.

\textbf{3.} Non : « descend sous $24$ » se traduit par \lstinline{temperature < 24}, test
strict que $24$ ne vérifie pas. Avec un seul seuil à $25\,^\circ\text{C}$, une température
oscillant entre $24{,}9$ et $25{,}1$ allumerait et éteindrait le chauffage à chaque mesure ;
les deux seuils ($24$ pour allumer, $26$ pour éteindre) laissent une zone où l'état
précédent est conservé, ce qui évite ces basculements incessants.
\end{corrige}
